% GENERATED FILE. Edit canonical lessons, then run npm run generate:books.
% canonical-source-hash: fnv1a64:7d13e374ffc0ef94
% canonical-lessons: SA-C368-sanka, SA-C368-tarkayati, SA-C368-asamsate, SA-C368-anumanam, SA-C368-kila

\chapter{Doubts and Hopes}
\label{ch:sa-doubts-and-hopes}
\hlchaptermodality{368}

\begin{chapteropening}
\textbf{By the end of this chapter:} \emph{I can say what I fear and what I hope will happen.}

Complete the last lesson of chapter 368: I can say what I fear and what I hope will happen.
\end{chapteropening}

\section[\texorpdfstring{\sk{शंका} (śaṅkā)}{शंका (śaṅkā)}]{\sk{शंका} (śaṅkā) — a misgiving; a fear}

\label{lesson:SA-C368-sanka}



\begin{quote}
Before the new one: say the Sanskrit for swiftly, then the Sanskrit for expects.
\end{quote}

\subsection*{You'll want to know: \sk{शंका}}
\textbf{\sk{शंका}} — \emph{śaṅkā} — \textquotedblleft{}a misgiving; a fear\textquotedblright{}.

\begin{quote}
\sk{यदि} \sk{शंका} \sk{अस्ति}, \sk{तर्हि} \sk{प्रश्नं} \sk{पृच्छतु।}
\end{quote}

\begin{quote}
\emph{yadi śaṅkā asti, tarhi praśnaṃ pṛcchatu.}
\end{quote}

\begin{quote}
If you have a misgiving, ask a question.
\end{quote}

A feminine noun.

\begin{cousinweb}
\textbf{\sk{शंका}} (\emph{śaṅkā}) comes from the root √śaṅk, \textquotedblleft{}to doubt, to fear\textquotedblright{}; MW gives \emph{śaṅkā} \textquotedblleft{}apprehension, fear, distrust\textquotedblright{} and \textquotedblleft{}doubt, hesitation\textquotedblright{} (MW).
\end{cousinweb}

\begin{grammarlens}[title={What you've built}]
One more word for saying what happens if.
\end{grammarlens}

\subsection*{Your turn}
\begin{itemize}
\raggedright
  \item \emph{Say it:} \emph{śaṅkā}
  \item \emph{Say it:} \emph{śaṅkā}, once more
  \item \emph{Say it:} \emph{yadi śaṅkā asti, tarhi praśnaṃ pṛcchatu.} — If you have a misgiving, ask a question.
  \item \emph{Recall:} say the Sanskrit for swiftly, then the Sanskrit for expects, then say \emph{śaṅkā} again
\end{itemize}

\begin{tcolorbox}[breakable,colback=teal!4,colframe=teal!35!black,title={Before you move on}]
What does \sk{शंका} mean? (\textquotedblleft{}A misgiving; a fear\textquotedblright{}.) Say it once more.
\end{tcolorbox}
\section[\texorpdfstring{\sk{तर्कयति} (tarkayati)}{तर्कयति (tarkayati)}]{\sk{तर्कयति} (tarkayati) — conjectures; guesses}

\label{lesson:SA-C368-tarkayati}



\begin{quote}
Before the new one: say the Sanskrit for expects, then the Sanskrit for a misgiving.
\end{quote}

\subsection*{You'll want to know: \sk{तर्कयति}}
\textbf{\sk{तर्कयति}} — \emph{tarkayati} — \textquotedblleft{}conjectures; guesses\textquotedblright{}.

\begin{quote}
\sk{यदि} \sk{दीपः} \sk{न} \sk{ज्वलति}, \sk{तर्हि} \sk{तर्कयामि} \sk{यत्} \sk{सः} \sk{गतः।}
\end{quote}

\begin{quote}
\emph{yadi dīpaḥ na jvalati, tarhi tarkayāmi yat saḥ gataḥ.}
\end{quote}

\begin{quote}
If no lamp is burning, I guess that he has gone.
\end{quote}

Past (he …): \sk{सः} \sk{कारणं} \sk{तर्कितवान्}, \emph{saḥ kāraṇaṃ tarkitavān}.

\begin{cousinweb}
\textbf{\sk{तर्कयति}} (\emph{tarkayati}) is built on \emph{tarka}, \textquotedblleft{}conjecture, reasoning\textquotedblright{}; MW gives \emph{√tark} \textquotedblleft{}to conjecture, to guess, to suspect, to infer\textquotedblright{} (MW).
\end{cousinweb}

\begin{grammarlens}[title={What you've built}]
One more word for saying what happens if.
\end{grammarlens}

\subsection*{Your turn}
\begin{itemize}
\raggedright
  \item \emph{Say it:} \emph{tarkayati}
  \item \emph{Say it:} \emph{tarkayati}, once more
  \item \emph{Say it:} \emph{yadi dīpaḥ na jvalati, tarhi tarkayāmi yat saḥ gataḥ.} — If no lamp is burning, I guess that he has gone.
  \item \emph{Recall:} say the Sanskrit for expects, then the Sanskrit for a misgiving, then say \emph{tarkayati} again
\end{itemize}

\begin{tcolorbox}[breakable,colback=teal!4,colframe=teal!35!black,title={Before you move on}]
What does \sk{तर्कयति} mean? (\textquotedblleft{}Conjectures; guesses\textquotedblright{}.) Say it once more.
\end{tcolorbox}
\section[\texorpdfstring{\sk{आशंसते} (āśaṃsate)}{आशंसते (āśaṃsate)}]{\sk{आशंसते} (āśaṃsate) — hopes for}

\label{lesson:SA-C368-asamsate}



\begin{quote}
Before the new one: say the Sanskrit for a misgiving, then the Sanskrit for conjectures.
\end{quote}

\subsection*{You'll want to know: \sk{आशंसते}}
\textbf{\sk{आशंसते}} — \emph{āśaṃsate} — \textquotedblleft{}hopes for\textquotedblright{}.

\begin{quote}
\sk{यदि} \sk{सः} \sk{प्रयत्नं} \sk{करोति}, \sk{तर्हि} \sk{सफलताम्} \sk{आशंसते।}
\end{quote}

\begin{quote}
\emph{yadi saḥ prayatnaṃ karoti, tarhi saphalatām āśaṃsate.}
\end{quote}

\begin{quote}
If he makes an effort, he hopes for success.
\end{quote}

Past (he …): \sk{सः} \sk{सफलताम्} \sk{आशंसितवान्}, \emph{saḥ saphalatām āśaṃsitavān}.

\begin{cousinweb}
\textbf{\sk{आशंसते}} (\emph{āśaṃsate}) is built on \emph{ā} and the root √śaṃs, \textquotedblleft{}to declare, to praise\textquotedblright{}; MW gives \emph{ā-√śaṃs} \textquotedblleft{}to hope for, to expect, to wish to attain\textquotedblright{} and \emph{ā-śaṃsita} \textquotedblleft{}hoped, expected\textquotedblright{} (MW).
\end{cousinweb}

\begin{grammarlens}[title={What you've built}]
One more word for saying what happens if.
\end{grammarlens}

\subsection*{Your turn}
\begin{itemize}
\raggedright
  \item \emph{Say it:} \emph{āśaṃsate}
  \item \emph{Say it:} \emph{āśaṃsate}, once more
  \item \emph{Say it:} \emph{yadi saḥ prayatnaṃ karoti, tarhi saphalatām āśaṃsate.} — If he makes an effort, he hopes for success.
  \item \emph{Recall:} say the Sanskrit for a misgiving, then the Sanskrit for conjectures, then say \emph{āśaṃsate} again
\end{itemize}

\begin{tcolorbox}[breakable,colback=teal!4,colframe=teal!35!black,title={Before you move on}]
What does \sk{आशंसते} mean? (\textquotedblleft{}Hopes for\textquotedblright{}.) Say it once more.
\end{tcolorbox}
\section[\texorpdfstring{\sk{अनुमानम्} (anumānam)}{अनुमानम् (anumānam)}]{\sk{अनुमानम्} (anumānam) — an inference; a guess}

\label{lesson:SA-C368-anumanam}



\begin{quote}
Before the new one: say the Sanskrit for conjectures, then the Sanskrit for hopes for.
\end{quote}

\subsection*{You'll want to know: \sk{अनुमानम्}}
\textbf{\sk{अनुमानम्}} — \emph{anumānam} — \textquotedblleft{}an inference; a guess\textquotedblright{}.

\begin{quote}
\sk{यदि} \sk{मम} \sk{अनुमानं} \sk{सत्यम्}, \sk{तर्हि} \sk{सः} \sk{श्वः} \sk{गमिष्यति।}
\end{quote}

\begin{quote}
\emph{yadi mama anumānaṃ satyam, tarhi saḥ śvaḥ gamiṣyati.}
\end{quote}

\begin{quote}
If my guess is right, he will go tomorrow.
\end{quote}

A neuter noun.

\begin{cousinweb}
\textbf{\sk{अनुमानम्}} (\emph{anumānam}) is built on \emph{anu}, \textquotedblleft{}after\textquotedblright{}, and the root √mā, \textquotedblleft{}to measure\textquotedblright{}; MW gives \emph{anu-māna} \textquotedblleft{}inference, consideration\textquotedblright{} and \textquotedblleft{}guess, conjecture\textquotedblright{} (MW).
\end{cousinweb}

\begin{grammarlens}[title={What you've built}]
One more word for saying what happens if.
\end{grammarlens}

\subsection*{Your turn}
\begin{itemize}
\raggedright
  \item \emph{Say it:} \emph{anumānam}
  \item \emph{Say it:} \emph{anumānam}, once more
  \item \emph{Say it:} \emph{yadi mama anumānaṃ satyam, tarhi saḥ śvaḥ gamiṣyati.} — If my guess is right, he will go tomorrow.
  \item \emph{Recall:} say the Sanskrit for conjectures, then the Sanskrit for hopes for, then say \emph{anumānam} again
\end{itemize}

\begin{tcolorbox}[breakable,colback=teal!4,colframe=teal!35!black,title={Before you move on}]
What does \sk{अनुमानम्} mean? (\textquotedblleft{}An inference; a guess\textquotedblright{}.) Say it once more.
\end{tcolorbox}
\section[\texorpdfstring{\sk{किल} (kila)}{किल (kila)}]{\sk{किल} (kila) — so they say; indeed}

\label{lesson:SA-C368-kila}



\begin{quote}
Before the new one: say the Sanskrit for hopes for, then the Sanskrit for an inference.
\end{quote}

\subsection*{You'll want to know: \sk{किल}}
\textbf{\sk{किल}} — \emph{kila} — \textquotedblleft{}so they say; indeed\textquotedblright{}.

\begin{quote}
\sk{सः} \sk{किल} \sk{रुग्णः}; \sk{यदि} \sk{एवम्}, \sk{तर्हि} \sk{अहं} \sk{गमिष्यामि।}
\end{quote}

\begin{quote}
\emph{saḥ kila rugṇaḥ; yadi evam, tarhi ahaṃ gamiṣyāmi.}
\end{quote}

\begin{quote}
He is ill, so they say; if so, I will go.
\end{quote}

\begin{cousinweb}
\textbf{\sk{किल}} (\emph{kila}) MW gives \emph{kila} as a particle of emphasis, \textquotedblleft{}indeed, verily\textquotedblright{}, and \textquotedblleft{}so said, so reported\textquotedblright{}; MW gives no derivation, so its origin is uncertain (MW).
\end{cousinweb}

\begin{grammarlens}[title={What you've built}]
One more word for saying what happens if.
\end{grammarlens}

\subsection*{Your turn}
\begin{itemize}
\raggedright
  \item \emph{Say it:} \emph{kila}
  \item \emph{Say it:} \emph{kila}, once more
  \item \emph{Say it:} \emph{saḥ kila rugṇaḥ; yadi evam, tarhi ahaṃ gamiṣyāmi.} — He is ill, so they say; if so, I will go.
  \item \emph{Recall:} say the Sanskrit for hopes for, then the Sanskrit for an inference, then say \emph{kila} again
\end{itemize}

\begin{tcolorbox}[breakable,colback=teal!4,colframe=teal!35!black,title={Before you move on}]
What does \sk{किल} mean? (\textquotedblleft{}So they say; indeed\textquotedblright{}.) Say it once more.
\end{tcolorbox}
